\chapter{Chú giải coverpoint của mô hình \FunctionalCoverage{}}
\label{app:cov-legend}

Phụ lục này giải nghĩa các coverpoint của mô hình \FunctionalCoverage{} trong
\cref{tab:functional-cov-matrix}. Các trường thuộc \CommandDescriptor{}, \ResponseDescriptor{},
mục \gls{dat} và \gls{csr} đã được mô tả tại \cref{app:csr,app:desc}.

\section*{Ý định phía \Host{}}

\begin{longtable}{L{4.2cm} L{8.4cm}}
  \toprule
  \textbf{Coverpoint} & \textbf{Ý nghĩa} \\
  \midrule
  \endhead
  \texttt{queue\_op}         & Thao tác phần mềm trên các queue port (đọc/ghi \CommandFifo{}, \TxRxFifo{}, \ResponseFifo{}). \\
  \texttt{bus\_enable}       & Bit bật bộ điều khiển trong \gls{csr}. \\
  \texttt{broadcast\_enable} & Bit cho phép phát broadcast header (\texttt{0x7E}) trước \AddressHeader{}. \\
  \texttt{hc\_abort}         & Bit \Host{} yêu cầu hủy giao dịch. \\
  \bottomrule
\end{longtable}

\section*{Quan sát trên bus (\texttt{i3c\_coverage})}

\begin{longtable}{L{4.2cm} L{8.4cm}}
  \toprule
  \textbf{Coverpoint} & \textbf{Ý nghĩa} \\
  \midrule
  \endhead
  \texttt{i3c}                     & Giao dịch tái tạo trên bus là I3C hay \iic{}. \\
  \texttt{i3c\_direct}             & CCC ở dạng Direct hay Broadcast. \\
  \texttt{CCC}                     & Mã lệnh CCC quan sát được (ENEC, DISEC, ENTDAA, \ldots). \\
  \texttt{addr\_class}             & Phân loại địa chỉ trên bus: broadcast (\texttt{0x7E}), I3C động, \iic{} tĩnh, hoặc dành riêng khác. \\
  \texttt{addr\_nack}              & Địa chỉ được ACK hay NACK. \\
  \texttt{start\_source}           & Điều kiện bắt đầu là START hay \RepeatedSTART{}. \\
  \texttt{bus\_op}                 & Chiều giao dịch trên bus (đọc hoặc ghi). \\
  \texttt{num\_data}               & Số byte dữ liệu của giao dịch. \\
  \texttt{trailing\_bytes}         & Số byte ở \gls{dword} cuối chưa đủ bốn (\texttt{num\_data} chia dư 4). \\
  \texttt{previous\_op}, \texttt{current\_op} & Chiều đọc/ghi của hai \PrivateTransfer{} liên tiếp nối qua \RepeatedSTART{}. \\
  \texttt{daa\_round\_outcome}     & Kết quả một vòng \gls{entdaa}: chấp nhận, địa chỉ bị từ chối, không có thiết bị, hoặc bị ngắt giữa chừng. \\
  \texttt{daa\_assigned\_addr\_class} & Vùng \DynamicAddress{} được gán (thấp, giữa, cao). \\
  \bottomrule
\end{longtable}

\section*{Tương quan đầu cuối (\texttt{i3c\_correlated\_coverage})}

\begin{longtable}{L{4.2cm} L{8.4cm}}
  \toprule
  \textbf{Coverpoint} & \textbf{Ý nghĩa} \\
  \midrule
  \endhead
  \texttt{target\_is\_i3c}       & \Target{} đích là I3C hay \iic{}. \\
  \texttt{nack\_phase}           & Trong \ImmediateDataTransfer{}, NACK xảy ra ở pha địa chỉ hay pha dữ liệu. \\
  \texttt{response\_len\_relation} & Quan hệ giữa độ dài dữ liệu thực tế và độ dài yêu cầu (đúng, thiếu, bằng 0, thiếu do abort, thiếu do overflow). \\
  \texttt{daa\_result}           & Kết quả tổng thể của \gls{entdaa} (gán hết, gán ít hơn yêu cầu, không thiết bị, bị từ chối, tràn, hủy). \\
  \texttt{preamble\_policy}      & Chính sách phát broadcast header, tổ hợp từ loại \Target{}, bit cho phép và header kỳ vọng so với header quan sát. \\
  \texttt{response\_addr\_phase} & \ResponseDescriptor{} ứng với pha địa chỉ của \Target{} hay của broadcast header. \\
  \texttt{address\_acked}        & Địa chỉ được ACK hay NACK, xét ở tầng tương quan. \\
  \texttt{response\_present}     & Giao dịch có sinh \ResponseDescriptor{} hay không. \\
  \texttt{final\_t\_bit}         & Giá trị \TBit{} cuối trong \PrivateRead{}: \Target{} kết thúc hay còn muốn truyền tiếp. \\
  \texttt{length\_outcome}       & Độ dài đọc thực tế so với yêu cầu: đúng, kết thúc sớm, hoặc vượt quá. \\
  \texttt{nack\_position}        & Vị trí byte bị NACK trong ghi \iic{}: đầu, giữa, cuối, hoặc không có. \\
  \texttt{short\_boundary}       & Vị trí cắt của giao dịch đọc kết thúc sớm so với ranh giới \gls{dword}. \\
  \texttt{protocol\_direction}   & Tổ hợp giao thức (I3C/\iic{}) và chiều (đọc/ghi). \\
  \texttt{abort\_cause}          & Nguyên nhân giao dịch dừng giữa chừng: \Host{} hủy, reset, hoặc giao thức kết thúc. \\
  \texttt{abort\_point}          & Pha trên bus khi giao dịch dừng (preamble, địa chỉ, TX, RX, CCC, \gls{entdaa}, hoặc \ResponseDescriptor{}). \\
  \texttt{recovery\_source}      & Nguồn khiến \gls{fsm} khôi phục xử lý lệnh, tương ứng ba nguyên nhân của \texttt{abort\_cause}. \\
  \texttt{command\_boundary}     & Cách một lệnh kết thúc: STOP, \RepeatedSTART{}, nối tiếp theo \texttt{toc}, idle liên tiếp, hoặc bị reset xóa. \\
  \texttt{stall\_type}           & Loại đình trệ: \TxFifo{} rỗng, \RxFifo{} đầy, \ResponseFifo{} đầy, hoặc chờ lệnh. \\
  \bottomrule
\end{longtable}
